%% SOURCE: experiments/010-kuzmin-tmi/scripts/positive_interaction_retrieval.py -- values read from experiments/010-kuzmin-tmi/results/positive_interaction_retrieval.csv; calibration from calibration_on_test_split.py, results/test_split_calibration.{csv,json}; prediction range from results/inference_1_validity.json
\section{Whether the model can pick positives\secstatus{todo}}\label{sec:positives}

Nominating positive interactions had not been tested. The earlier retrieval
measurement used the Kuzmin 2018 call, which is one-sided by construction:
``we used an established interaction magnitude cut-off for digenic interactions
$(p < 0.05, |\varepsilon| > 0.08)$ and trigenic interactions
$(p < 0.05, \tau < -0.08)$'', with positives excluded because
``We focused exclusively on the analysis of deleterious negative trigenic
interactions.'' The symmetric form belongs to Kuzmin 2020 and the 2021 protocol,
which do score positives, and it is what makes the question askable on the same
records.

\subsection{Positives are well represented in the labels\secstatus{todo}}

Of the 376{,}732 records in this build, 20{,}426 have $\tau > +0.08$ and 1{,}773
have $\tau > +0.20$, against a label standard deviation of $0.0633$ and a maximum
of $1.128$. The model trains on them.

\subsection{Retrieval, and the margin is wider than on negatives\secstatus{todo}}

Precision at $K = 100$ on the held-out test split, 37{,}673 records.

\begin{table}[H]
\centering
\begin{tabular}{lrrr}
\toprule
Model & P@100 & enrichment & average precision \\
\midrule
\multicolumn{4}{l}{\emph{Kuzmin 2020 call, $\tau > +0.08$ and $p < 0.05$, base rate 0.526\%}} \\
B1 additive ridge & 0.030 & 5.7$\times$  & 0.0110 \\
B5 nonlinear MLP  & 0.040 & 7.6$\times$  & 0.0166 \\
CGT ensemble      & 0.100 & 19.0$\times$ & 0.0352 \\
\midrule
\multicolumn{4}{l}{\emph{magnitude only, $\tau > +0.20$, base rate 0.422\%}} \\
B1 additive ridge & 0.040 & 9.5$\times$  & 0.0109 \\
B5 nonlinear MLP  & 0.100 & 23.7$\times$ & 0.0377 \\
CGT ensemble      & 0.210 & 49.8$\times$ & 0.1110 \\
\bottomrule
\end{tabular}
\caption{Retrieving published positive trigenic interactions on the
random-over-records test split. The transformer's average precision is
$3.2$ times the additive null's at the published-style call and $10.2$ times at
$\tau > +0.20$. On the negative side the same ratio was $1.9$.}
\end{table}

The margin over a model that cannot represent interaction is proportionally
larger on positives than on negatives, which is the opposite of what the
one-sided published call would suggest.

\subsection{Large predictions are calibrated on this split\secstatus{todo}}

The previous panel's collapsed doubles scored near $0.71$, more than ten label
standard deviations above the mean, and that inflation was the signature of the
bug. A top prediction on \texttt{inference\_1} is $+0.690$, so the same check is
owed here.

Binning the held-out test predictions against actual $\tau$: the 18 records
predicted above $+0.20$ have mean actual $\tau$ of $+0.312$, and the 5 predicted
above $+0.30$ have mean actual $\tau$ of $+0.802$. Over the whole split a
regression of actual on predicted has slope $0.877$. Test-split predictions
themselves span $-0.576$ to $+0.631$, so a prediction near $+0.65$ is inside the
range the model demonstrably emits on labeled data rather than outside it.

Two qualifications belong with those numbers, and both were missing from an
earlier draft.

\pfstep{Magnitude is not a call} Of the 18 records predicted above $+0.20$, 12
exceed $\tau > +0.08$ in magnitude but only \textbf{2} also clear $p < 0.05$, and
of the 5 predicted above $+0.30$ all five exceed the magnitude while
\textbf{none} clears the significance half. The enrichment is real either way,
$66.7$ percent against a magnitude base rate of $5.27$ percent and $11.1$ percent
against a called base rate of $0.53$ percent. But the Kuzmin 2020 call is the
conjunction, it is what the retrieval table above scores, and a magnitude-only
share must not be reported as a real call.

\pfstep{A whole-split slope is not the operating point} The $0.877$ is computed
over all 37{,}673 test records and is carried by the bulk near zero. Inside the
predicted head the same regression gives $1.59$ at the top 100 and $1.19$ at the
top 500, so on this split the model under-predicts where it ranks highest rather
than inflating. That is the opposite of what the same checkpoints do on the
\texttt{inference\_4} space, where the held-out slope is $0.24$.
Section~\ref{sec:panel20} reconciles the two, and the resolution decides how a
predicted $\tau$ may be read.

\subsection{Confidence now tracks presence of evidence\secstatus{todo}}

Over the 200 most positive \texttt{inference\_1} triples, the Spearman
correlation between prediction and the least-supported gene's trigenic record
count is $+0.267$ at $p = 1.3\times10^{-4}$, and between prediction and the mean
record count $+0.279$. Thirty-one of the 200 contain a gene with no record, not
the whole list. The direction is the reverse of the previous panel's.

The panel genes back this at the label level. \texttt{YER059W} appears in 1{,}097
training triples of which $16.3$ percent exceed $+0.08$, a $3.0$ times
enrichment over the $5.4$ percent base rate. \texttt{YKR028W} reaches $14.7$
percent. The model's preference is visible in the labels it was trained on.

\subsection{What none of this covers\secstatus{todo}}

Every number above comes from the split that is random over records, where a
Kuzmin query double appears in training and test at once. On a query-pair
disjoint split the additive null falls from $0.400$ to $0.127 \pm 0.033$ test
Pearson and the nonlinear baseline drops below it. The transformer has never been
refit there. Only 22 of the 420 query doubles have both members inside the
526-gene space, so \texttt{inference\_1} is very nearly query-pair disjoint by
construction and its nominations live in exactly the regime that has not been
measured.

The positive side of that split is underpowered rather than clearly collapsed.
It holds 51 published positive calls in 37{,}680 records, a base rate of $0.135$
percent, and the additive ridge retrieves 3 of them in its top 100. The
enrichment ratio that produces, $22.2$ times, is not comparable to the random
split's and rests on three records. Negative retrieval on the same split falls to
$1.8$ times, which is close to chance, but the positive comparison cannot yet be
called either way.
